\documentclass[10pt,a4paper]{article}
% =====================================================================
%  État d'avancement — version compacte (≤ 7 pages), lisible et expliquée
%  Rédigé le 2026-10-02. Chaque chiffre porte sa date et son artefact.
%  Compilation : xelatex
% =====================================================================
\usepackage{fontspec}
\setmainfont{TeX Gyre Pagella}
\setsansfont{Lato}
\setmonofont{DejaVu Sans Mono}[Scale=0.84]
\usepackage[french]{babel}
\frenchsetup{AutoSpacePunctuation=true}
\usepackage[margin=1.55cm,top=1.7cm,bottom=1.5cm,headheight=14pt]{geometry}
\usepackage{amsmath,amssymb}
\usepackage{booktabs,colortbl,xcolor}
\usepackage{tcolorbox}
\tcbuselibrary{skins,breakable}
\usepackage{tikz}
\usetikzlibrary{positioning,arrows.meta,shapes.geometric,calc,fit,backgrounds,decorations.pathreplacing}
\usepackage{pgfplots}
\pgfplotsset{compat=1.18}
\usepackage{fontawesome5}
\usepackage{titlesec}
\usepackage{fancyhdr}
\usepackage[shortlabels]{enumitem}
\usepackage{multicol}
\usepackage{caption}
\usepackage{hyperref}
\hypersetup{colorlinks=true,linkcolor=primary,urlcolor=primary!80!black,pdftitle={État d'avancement au 2 octobre 2026},pdfauthor={Mohamed Salhi}}

\providecommand{\NoAutoSpacing}{}
\newcommand{\code}[1]{{\NoAutoSpacing\ttfamily\detokenize{#1}}}
\newcommand{\pct}[1]{#1\,\%}

\definecolor{primary}{HTML}{1B3A5C}
\definecolor{accent}{HTML}{0E7490}
\definecolor{ok}{HTML}{15803D}
\definecolor{warn}{HTML}{B45309}
\definecolor{fail}{HTML}{B91C1C}
\definecolor{todo}{HTML}{64748B}
\definecolor{ink}{HTML}{111827}
\definecolor{paper}{HTML}{F7F5F1}

\newcommand{\OK}{\textcolor{ok}{\faCheckCircle}}
\newcommand{\KO}{\textcolor{fail}{\faTimesCircle}}
\newcommand{\PART}{\textcolor{warn}{\faClock}}
\newcommand{\TODO}{\textcolor{todo}{\faHourglassHalf}}

\captionsetup{font={footnotesize,sf},labelfont={bf,color=primary},skip=3pt}
\captionsetup[table]{skip=3pt}

\tikzset{
  boite/.style={draw=primary!55, fill=primary!4, rounded corners=2pt,
    align=center, inner sep=4pt, font=\scriptsize\sffamily},
  boiteok/.style={boite, draw=ok!70, fill=ok!7},
  boiteaccent/.style={boite, draw=accent!80, fill=accent!7},
  boitetodo/.style={boite, draw=todo!70, fill=todo!7, dashed},
  fleche/.style={-{Stealth[length=2.2mm]}, draw=primary!70, thick},
}

\newtcolorbox{clair}[1][]{colback=accent!6, colframe=accent!70, boxrule=0.5pt,
  left=5pt,right=5pt,top=4pt,bottom=4pt, fonttitle=\bfseries\sffamily\small,
  title={\faLightbulb~En clair}, #1}
\newtcolorbox{atrancher}[1][]{colback=warn!6, colframe=warn!75, boxrule=0.5pt,
  left=5pt,right=5pt,top=4pt,bottom=4pt, fonttitle=\bfseries\sffamily\small,
  title={\faExclamationTriangle~À trancher}, #1}
\newtcolorbox{limite}[1][]{colback=fail!5, colframe=fail!65, boxrule=0.5pt,
  left=5pt,right=5pt,top=4pt,bottom=4pt, fonttitle=\bfseries\sffamily\small,
  title={\faBan~Ce que le système n'a pas le droit de dire}, #1}

\titleformat{\section}[block]
  {\normalfont\sffamily\large\bfseries\color{primary}}{\thesection.}{0.5em}{}
  [{\vspace{-6pt}\color{primary!25}\rule{\linewidth}{0.8pt}}]
\titleformat{\subsection}[block]
  {\normalfont\sffamily\normalsize\bfseries\color{primary!88}}{\thesubsection}{0.45em}{}
\titlespacing*{\section}{0pt}{8pt}{4pt}
\titlespacing*{\subsection}{0pt}{6pt}{2pt}
\setlist[itemize]{leftmargin=1.2em,itemsep=1pt,topsep=2pt}
\renewcommand{\arraystretch}{1.08}
\setlength{\tabcolsep}{3pt}
\pagestyle{fancy}
\fancyhf{}
\fancyhead[L]{\scriptsize\sffamily\color{primary}État d'avancement — 2 octobre 2026}
\fancyhead[R]{\scriptsize\sffamily\color{primary}Audit IoT · reconnaissance d'objets connectés}
\fancyfoot[C]{\scriptsize\sffamily\color{todo}\thepage}
\renewcommand{\headrulewidth}{0.3pt}
\renewcommand{\headrule}{\hbox to\headwidth{\color{primary!30}\leaders\hrule height \headrulewidth\hfill}}

\begin{document}
\sloppy
\emergencystretch=2em

% =====================================================================
%  BANDEAU DE TITRE (compact)
% =====================================================================
\begin{tikzpicture}[remember picture,overlay]
  \fill[primary] (current page.north west) rectangle ([yshift=-2.35cm]current page.north east);
  \fill[accent] ([yshift=-2.35cm]current page.north west) rectangle ([yshift=-2.49cm]current page.north east);
\end{tikzpicture}
\vspace*{0.15cm}
{\color{white}\sffamily
\noindent{\LARGE\bfseries État d'avancement — 2 octobre 2026}\\[1mm]
\noindent{\small Reconnaissance d'objets connectés pour l'audit de sécurité IoT · ce qui est atteint, mesuré et prouvé}\\[0.5mm]
\noindent{\footnotesize\color{white!85}Mohamed Salhi — Polytech Lille / ENSA Kénitra · Tuteur : Thomas Vantroys · Horizon 15/09 → 15/12/2026}}
\vspace{1.35cm}

{\footnotesize\sffamily
\noindent\textbf{Comment lire.} Ce document ne raconte pas ce qui est prévu : il raconte ce qui a
été \textbf{mesuré}, avec la date et l'endroit où la preuve est rangée. Statuts : \OK{} fait et
prouvé · \PART{} commencé mais incomplet · \TODO{} pas commencé. Les mesures vont jusqu'au
\textbf{29/09/2026} ; le document est écrit le \textbf{02/10/2026}. Les sigles et les règles de
comptage sont expliqués au §2, et \textbf{chaque chiffre est accompagné de ce qu'il veut dire}.}

\section{Ce qui est atteint, en quatre phrases}
\begin{center}
\renewcommand{\arraystretch}{1.0}
\begin{tabular}{@{}p{0.48\linewidth}p{0.48\linewidth}@{}}
\begin{tcolorbox}[enhanced,colback=ok!5,colframe=ok!60,boxrule=0.5pt,arc=2pt,
  left=5pt,right=5pt,top=4pt,bottom=4pt]
{\sffamily\bfseries\color{ok}1. Les machines} {\footnotesize\sffamily Trois ordinateurs de
l'école font tourner le modèle d'IA. Tout se lance, se mesure et s'arrête par des scripts. Un test
de charge relancé 13 jours plus tard donne \textbf{les mêmes chiffres à 1\,\% près}. Après une
coupure de courant, le service revient \textbf{tout seul en 75 secondes} et répond. Et internet
coupé, l'IA continue : c'est vérifié pour de vrai, sur la machine.}
\end{tcolorbox}&
\begin{tcolorbox}[enhanced,colback=warn!5,colframe=warn!60,boxrule=0.5pt,arc=2pt,
  left=5pt,right=5pt,top=4pt,bottom=4pt]
{\sffamily\bfseries\color{warn}2. La reconnaissance} {\footnotesize\sffamily Sur 48 objets
photographiés, l'outil en identifie complètement \textbf{10} (la cible était 34 sur 48). La cause
est trouvée : quand la marque est \textbf{lisible} sur la photo, il la lit bien (9 objets sur 17
parfaitement identifiés, jamais une marque fausse) ; quand elle ne l'est pas, c'est presque
toujours raté (1 sur 31, et 12 marques fausses). En résumé : il \textbf{lit les étiquettes}, il ne
\textbf{reconnaît pas les formes}.}
\end{tcolorbox}\\[1.5mm]
\begin{tcolorbox}[enhanced,colback=ok!5,colframe=ok!60,boxrule=0.5pt,arc=2pt,
  left=5pt,right=5pt,top=4pt,bottom=4pt]
{\sffamily\bfseries\color{ok}3. Les failles} {\footnotesize\sffamily La base officielle des
failles (NVD, l'agence américaine qui les recense) est \textbf{copiée en local} : 121\,719 failles
connues. Une recherche répond en \textbf{0,27 milliseconde} là où le contrat en autorise 50 :
\textbf{près de 200 fois plus rapide} que le minimum exigé. Mais un choix de filtre fait que, pour
les marques présentes dans nos photos, \textbf{10 failles seulement sur 11\,714} sont conservées :
il faut décider de l'élargir.}
\end{tcolorbox}&
\begin{tcolorbox}[enhanced,colback=todo!5,colframe=todo!60,boxrule=0.5pt,arc=2pt,
  left=5pt,right=5pt,top=4pt,bottom=4pt]
{\sffamily\bfseries\color{todo}4. Ce qui reste à faire} {\footnotesize\sffamily Le rapport PDF
signé (avec un journal qu'on ne peut pas modifier), la page web de démonstration, le calcul de
confiance croisée (accord entre photos + fiche officielle existante) et l'index de recherche rapide
des produits. Les jalons G3 à G7 — expliqués au §2 — sont prévus de \textbf{novembre à décembre}.}
\end{tcolorbox}
\end{tabular}
\end{center}

\begin{clair}
\textbf{Les trois phrases à retenir.} (1)~L'infrastructure n'est plus un risque : elle est finie,
mesurée et rejouable. (2)~Le cœur du sujet est mesuré honnêtement bas, et son obstacle est
identifié : il faut des photos où la marque est \textbf{lisible}. (3)~La recherche de failles
fonctionne, mais un filtre trop strict ferait afficher « aucune faille » pour des objets qui en
ont : \textbf{c'est la décision la plus urgente.}
\end{clair}

% =====================================================================
\section{Repères de lecture (à garder sous les yeux)}
% =====================================================================
{\footnotesize\sffamily\emergencystretch=3em
\begin{multicols}{2}
\begin{itemize}
\item \textbf{Objet / image.} Un \emph{objet} (48 au total) a plusieurs photos : 144 au total.
Un score « par objet » consolide toutes ses photos ; un score « par image » compte chaque photo.
\item \textbf{Marque / modèle.} La \emph{marque} est le fabricant (« TP-Link ») ; le \emph{modèle}
est la référence exacte (« TL-WR841N »). « Identifier » = donner les deux.
\item \textbf{Catégorie.} Le \emph{type} d'objet : routeur, caméra, prise, ampoule, enceinte.
C'est la question la plus facile ; trouver le modèle exact est la plus dure.
\item \textbf{IC95 (« les crochets »).} La fourchette d'incertitude à 95\,\% : « 10/48 =
20,8\,\% [11,7 ; 34,3] » se lit « on a mesuré 20,8\,\%, la vraie valeur est probablement entre
11,7\,\% et 34,3\,\% ». Deux fourchettes qui ne se recouvrent pas = écart réel, pas du hasard.
Tous les crochets de ce document se lisent ainsi.
\item \textbf{Barème (A / P) — et les deux barèmes du projet.} Un \emph{barème} est la règle de
comparaison appliquée avant de compter : \emph{A} = stricte (« DIR-600 » $\neq$ « dir 600 ») ;
\emph{P} = tirets et espaces ignorés. Le modèle écrit souvent la \emph{bonne} référence avec un
tiret différent : les deux résultats se publient. Le projet a deux barèmes : \code{metrics} (le
barème interne, qui compte par objet) et \code{audit-grading} (un barème externe figé, recopié
dans le dépôt, qui compte par photo).
\item \textbf{Règles de comptage.} Une \emph{abstention} = le modèle répond « je ne sais pas ».
La \emph{règle publiée} la compte comme une erreur ; la \emph{règle des réponses données} ne note
que les objets où le modèle a osé répondre. \emph{Plafond} : borne haute optimiste, où une seule
bonne photo par objet suffit.
\item \textbf{« Juste » / « faux ».} \emph{Juste} = la réponse est identique à l'étiquetage humain
de référence ; \emph{fausse} = différente.
\item \textbf{M1.} La mesure principale : \emph{marque et modèle exacts} pour un objet, après
consolidation de ses photos. Cible : 70\,\%. « \emph{v4 figée} » désigne l'exécution de référence
dont les réponses servent partout dans ce document.
\item \textbf{Sigles.} \emph{VLM} : le modèle de vision qui regarde la photo. \emph{CPE} : le code
officiel d'un modèle de produit. \emph{CVE} : le nom public d'une faille. \emph{NVD} : la base
officielle qui relie failles et produits. \emph{OCR} : la lecture automatique du texte imprimé.
\item \textbf{Strates.} Les familles de photos : \emph{officiel} 3 (visuels constructeur — le
modèle a pu les voir pendant son apprentissage), \emph{avis} 129, \emph{forums/démontages} 12
(angles inhabituels, test le plus honnête).

\item \textbf{Portes G0–G7.} Les huit étapes de validation du calendrier : une « porte » ne se
ferme qu'en \emph{montrant la preuve}, pas en promettant. \textbf{G0} socle + corpus (fermée le
24/09) · \textbf{G1} base de failles rejouée + 60 cas vérifiés (dû 19/10) · \textbf{G2} recherche
produits + 200 cas (02/11) · \textbf{G3} chaîne complète photo → PDF (23/11) · \textbf{G4} journal
scellé + signature + tests de falsification (30/11) · \textbf{G5} interface web (07/12) ·
\textbf{G6} tous les chiffres régénérés (14/12) · \textbf{G7} gel final (15/12).
\item \textbf{« Niveau × barème ».} Règle d'écriture du projet : un chiffre doit toujours dire
\emph{comment} il a été compté (objet ou image ; barème A ou P ; règle de comptage), sinon il n'est
pas citable. C'est ce que ce document applique partout.
\end{itemize}
\end{multicols}}

\subsection{Les chiffres-phares, et ce qu'ils veulent dire}
\begin{center}
\footnotesize\sffamily
\begin{tabular}{@{}p{0.235\linewidth}p{0.30\linewidth}p{0.40\linewidth}@{}}
\toprule
\textbf{Ce qui a été mesuré} & \textbf{Résultat} & \textbf{Ce que ça veut dire}\\
\midrule
Identification complète & 10 objets sur 48 & Loin de la cible de 34/48 — mais la cause est connue (voir §3).\\
Marque lisible dans la photo & 9 objets sur 17 identifiés ; 0 marque inventée & Là où il y a du texte, il lit bien.\\
Marque illisible & 1 objet sur 31 & Sans texte, il ne devine pas la marque — c'est mesuré, pas supposé.\\
Relancer le même test & 144 photos sur 144 donnent la même réponse & Le système est stable : les chiffres sont fiables.\\
Temps par photo & médiane 2,4 s ; 95\,\% des photos sous 3,8 s & La cible contractuelle est 10 s : tenue avec marge.\\
Base de failles locale & 121\,719 failles · 324\,470 règles de correspondance & L'équivalent de plusieurs années du NVD, utilisable hors-ligne.\\
Recherche d'une faille & 0,27 ms (cible : 50 ms) & Le résultat sort instantanément, sans réseau.\\
Charge des 3 machines & 30,2 requêtes/s par machine au point retenu (16 requêtes simultanées) & Même valeur retrouvée à 13 jours d'écart ; au-delà, la latence monte plus vite que le débit.\\
Tests automatiques & 60 sur 60 passent & Les fondations ne cassent pas quand on modifie le code.\\
\bottomrule
\end{tabular}
\end{center}

% =====================================================================
\section{Le projet, et la limite qu'il ne peut pas franchir}
% =====================================================================
À partir d'une photo d'un seul objet connecté, l'outil doit deviner marque et modèle, retrouver la
fiche officielle du produit, puis lister les failles connues — sans jamais se connecter à internet
pendant un audit. Une photo montre le boîtier, jamais le logiciel interne : le système pourra donc
dire « ce modèle a connu des failles », jamais « cet exemplaire est vulnérable ». C'est une règle
inscrite dans le logiciel, pas une précaution de rédaction.

\begin{center}
\begin{tikzpicture}[node distance=2.6mm and 2.6mm, font=\scriptsize\sffamily]
\node[boiteok] (a) {\faCamera\ \textbf{1. Regarder}\\ photo → marque, modèle};
\node[boiteok,right=of a] (b) {\faFileSignature\ \textbf{2. Consolider}\\ accord entre photos + fiche};
\node[boiteaccent,right=of b] (c) {\faSearch\ \textbf{3. Résoudre}\\ fiche produit CPE + score};
\node[boiteok,right=of c] (d) {\faDatabase\ \textbf{4. Chercher}\\ failles connues, « niveau produit »};
\node[boitetodo,right=of d] (e) {\faFilePdf\ \textbf{5. Rapporter}\\ journal scellé + PDF \TODO};
\draw[fleche] (a)--(b); \draw[fleche] (b)--(c); \draw[fleche] (c)--(d);
\draw[fleche,dashed] (d)--(e);
\end{tikzpicture}
\captionof{figure}{La chaîne complète. La flèche pointillée est la seule étape non commencée à ce
jour. Chaque étape échange un formulaire fixe, ce qui permet de la tester seule.}\label{fig:chaine}
\end{center}

\begin{limite}
Une photo ne montre jamais la version du logiciel interne. Le vocabulaire du rapport est donc
verrouillé : on parle du \emph{modèle}, pas de l'exemplaire. Sont interdits, et testés comme tels :
« ces failles concernent votre exemplaire » · « aucune faille trouvée, donc l'objet est sûr » (rien
de répertorié $\neq$ rien qui existe) · « le système détecte les failles » ou une note de risque
agrégée · un pourcentage sans son effectif ni son intervalle.
\end{limite}

% =====================================================================
\section{La reconnaissance : ce qui marche, ce qui ne marche pas, et pourquoi}
% =====================================================================
\subsection{Le modèle choisi, et pourquoi}
Trois modèles d'IA ont regardé \textbf{les mêmes 144 photos}, avec la même consigne, sans hasard
(température 0), et notés contre la vérité :
\begin{center}
\footnotesize\sffamily
\begin{tabular}{@{}p{0.30\linewidth}p{0.26\linewidth}p{0.38\linewidth}@{}}
\toprule
\textbf{Modèle} & \textbf{Bonnes réponses /144} & \textbf{Avantage ou défaut décisif}\\
\midrule
\textbf{Qwen2.5-VL-7B} \emph{(retenu)} & \textbf{79} [0,467 ; 0,628] & Meilleur score ; \textbf{aucune} réponse hors format.\\
Llama-3.2-11B-Vision & 72 [0,419 ; 0,581] & Écarté pour la fiabilité : \textbf{21 réponses sur 144 hors format}, et 54 fois il affirme avoir lu une étiquette alors que rien n'est lisible.\\
Qwen3-VL-8B & 53 [0,294 ; 0,449] & Plus récent mais nettement moins bon ; l'écart entre les intervalles est réel.\\
\bottomrule
\end{tabular}

{\scriptsize\sffamily Lecture des crochets : c'est la fourchette à 95\,\% (voir §2) — « 79 sur
144 = 54,9\,\%, entre 46,7\,\% et 62,8\,\% ».}
\end{center}
Le choix n'est donc pas « le plus récent » : c'est le plus juste sur ce corpus \emph{et} celui qui
respecte toujours le format demandé — indispensable pour un outil automatique.

\subsection{Trois niveaux de difficulté, très différents}
La même photo peut être bien \emph{classée} (c'est une caméra) et mal \emph{identifiée} (une marque
fausse). Voici les trois niveaux, par photo :
\begin{center}
\footnotesize\sffamily
\begin{tabular}{@{}p{0.24\linewidth}p{0.22\linewidth}p{0.46\linewidth}@{}}
\toprule
\textbf{Niveau} & \textbf{Résultat} & \textbf{Ce que ça veut dire}\\
\midrule
Type d'objet (catégorie) & 122 / 144 (84,7\,\%) & Distinguer routeur, caméra, prise, ampoule, enceinte : le plus facile.\\
Modèle du produit & 79 / 144 (54,9\,\% en P) & Le nom exact (« TL-WR841N ») est plus dur que le type.\\
Marque seule & 46 / 144 (31,9\,\%) & Même la marque est rarement juste — d'où le diagnostic ci-dessous.\\
Marque \emph{et} modèle & 19 / 144 (13,2\,\%) & Le niveau le plus dur, celui que le produit doit atteindre.\\
\bottomrule
\end{tabular}

{\footnotesize\sffamily Ces quatre lignes comptent \emph{photo par photo} (barème \code{audit-grading}) ; le score par objet (M1, §4.3) est plus dur : il exige la bonne réponse \emph{après} avoir mis d'accord les photos du même objet.}
\end{center}

\subsection{Le score principal : identifier un objet entier (marque \emph{et} modèle)}
Un objet a plusieurs photos : la mesure principale du projet, \textbf{M1}, exige la bonne marque
\emph{et} le bon modèle pour l'objet, après avoir mis ses photos d'accord entre elles. Trois façons
de compter existent — et les trois se publient, parce qu'elles mesurent trois choses :

\begin{center}
\begin{minipage}[t]{0.47\linewidth}
\footnotesize\sffamily
\begin{tabular}{@{}p{0.30\linewidth}p{0.22\linewidth}p{0.40\linewidth}@{}}
\toprule
\textbf{Manière de compter} & \textbf{Résultat} & \textbf{Ce qu'elle mesure}\\
\midrule
Règle publiée (l'abstention compte) & \textbf{10/48 = 20,8\,\%} & Un « je ne sais pas » est compté comme une erreur.\\
Réponses données seulement & 15/48 = 31,2\,\% & On ne note que les objets où le modèle a osé répondre.\\
Plafond optimiste & 17/48 = 35,4\,\% & Même avec une seule bonne photo par objet, impossible d'aller au-delà.\\
\bottomrule
\end{tabular}
\end{minipage}\hfill
\begin{minipage}[t]{0.49\linewidth}
\begin{tikzpicture}
\begin{axis}[width=\linewidth,height=3.8cm,
  ybar,bar width=10pt,ymin=0,ymax=0.8,
  symbolic x coords={{règle publiée},{réponses},{plafond}},
  xtick=data, ylabel={M1 (par objet)}, nodes near coords,
  enlarge x limits=0.28, grid=major, grid style={primary!12},
  axis line style={primary!50}, tick style={primary!50},
  every node near coord/.append style={font=\scriptsize\sffamily\bfseries}]
\addplot[fill=primary!55,draw=primary!80] coordinates {
  ({règle publiée},0.208) ({réponses},0.312) (plafond,0.354)};
\draw[fail,thick,dashed] (axis cs:{règle publiée},0.70)--(axis cs:plafond,0.70)
  node[pos=1,right,font=\scriptsize\sffamily,color=fail]{cible 0,70};
\end{axis}
\end{tikzpicture}
\captionof{figure}{Les trois façons de compter le score M1 (par objet, exécution v4 figée, 144 photos / 48 objets).
La barre « plafond » montre qu'\textbf{aucun mode de comptage} ne rapproche de la cible : le
problème est dans ce que le modèle perçoit, pas dans la façon de compter.}
\end{minipage}
\end{center}

\begin{clair}
Pourquoi trois colonnes au lieu d'un seul chiffre ? Parce que le projet \textbf{ordonne} au modèle
de répondre « je ne sais pas » quand la marque n'est pas visible. La règle publiée compte cette
honnêteté comme une erreur ; la deuxième ne juge que les réponses données. Publier une seule des
deux, ce serait choisir une politique sans le dire. Et le plafond (17/48) fixe la limite
optimiste : même en offrant au système une photo parfaite par objet, on reste à 35\,\%.
\end{clair}

\subsection{Le résultat qui explique tout : la marque se \emph{lit}, elle ne se \emph{reconnaît} pas}
\begin{center}
\begin{tikzpicture}
\begin{axis}[width=0.60\linewidth,height=4.1cm,
  ybar,bar width=7pt,ymin=0,ymax=1.05,
  symbolic x coords={{marque (photo)},{objet exact},{catégorie (photo)}},
  xtick=data, ylabel={exactitude},
  legend style={font=\scriptsize\sffamily,at={(0.5,1.02)},anchor=south,legend columns=2},
  enlarge x limits=0.22, grid=major, grid style={primary!12},
  axis line style={primary!50}, tick style={primary!50},
  every node near coord/.append style={font=\scriptsize\sffamily}]
\addplot[fill=ok!70,draw=ok!80] coordinates {({marque (photo)},0.766) ({objet exact},0.529) ({catégorie (photo)},1.000)};
\addplot[fill=fail!60,draw=fail!80] coordinates {({marque (photo)},0.103) ({objet exact},0.032) ({catégorie (photo)},0.773)};
\legend{texte lisible sur la photo, texte illisible}
\end{axis}
\end{tikzpicture}
\captionof{figure}{Comment lire ce graphique : chaque paire de barres compare les photos où le
texte est lisible (vert) à celles où il ne l'est pas (rouge). Quand il y a du texte,
\textbf{77\,\% des photos ont la bonne marque} et 53\,\% des objets sont parfaitement identifiés ;
sans texte, ces taux tombent à \textbf{10\,\%} et 3\,\% — la marque juste est
\textbf{7,4 fois plus rare}.}\label{fig:canal}
\end{center}

Répartition complète des 48 objets (règle publiée) : \textbf{10 exacts} (marque et modèle) ·
\textbf{9} marque juste mais modèle faux · \textbf{17} catégorie juste mais marque non trouvée ·
\textbf{12} franchement faux. Les 12 cas « franchement faux » viennent \emph{tous} de photos sans
texte lisible. Quand le texte est lisible, le système ne donne \emph{jamais} une marque fausse :
soit il trouve la bonne, soit il se tait. Sans texte, 12 objets sur 31 reçoivent une marque
fausse. Ce n'est pas une dégradation progressive, c'est \emph{un canal qui s'éteint}. Exemple
parlant : sur l'enceinte \code{amazon\_\_echo\_dot}, la marque « amazon » est juste \textbf{12 fois
sur 12} quand le texte est lisible, et la réponse est vide \textbf{25 fois sur 29} quand il ne
l'est pas. Autrement dit : le système \textbf{lit} les étiquettes ; il ne \textbf{reconnaît} pas
les formes. Le levier n'est donc pas le prompt, mais \textbf{le corpus de photos} : il faut
photographier la marque (face, étiquette, sérigraphie).

Deux observations utiles, mesurées. (1)~Sur les 12 photos de \emph{démontage} (angles
inhabituels), \textbf{0 objet est identifié exactement} — normal : il n'y a rien à lire.
(2)~Les 3 \emph{visuels officiels} du constructeur ne sont pas mieux lotis sur l'identification
exacte (0 sur 3), mais leur catégorie est juste 3 fois sur 3 : c'est plutôt une bonne nouvelle,
car rien ne montre que le modèle \emph{récite} des images qu'il aurait vues pendant son
apprentissage.

\subsection{Trois pistes essayées pour réparer la marque — et pourquoi aucune ne marche}
\begin{itemize}
\item \textbf{Chercher la marque dans le dictionnaire officiel des produits (CPE) : 0 réussite sur
22 essais.} L'idée était : quand le modèle ne dit rien, chercher « Echo Dot » dans les 1,8 million
de fiches produits officielles. Résultat : ce dictionnaire est rempli de logiciels, pas d'objets
grand public — il propose \code{dnnsoftware} pour « Echo Dot » ou \code{minio} pour « HomePod
mini ». Sur 22 essais, \textbf{17 propositions sont fausses}. C'est aussi une explication du
mauvais score de la « résolution de fiche produit » (§8, AC10) : ce que ce dictionnaire ne décrit
pas, le système ne peut pas le retrouver.
\item \textbf{Changer la consigne pour demander la marque (essai « v5 ») : +1 objet seulement}
(11 au lieu de 10 sur 48), avec des fourchettes d'incertitude qui se chevauchent — l'écart peut
donc être du hasard ; et si l'on ne compte que les réponses données, c'est exactement pareil
(15/48 des deux côtés).
\item \textbf{Combiner deux consignes ne peut pas aider.} Vérifié : les deux versions laissent
\emph{les mêmes photos} sans réponse. C'est la photo qui décide, pas la formulation de la question.
\end{itemize}

\subsection{Le diagnostic chiffré : le prompt ne peut pas sauver la mise}
On a repris les photos ratées une par une — 65 photos sur 144 n'ont pas donné la bonne réponse —
et demandé à \textbf{deux lecteurs de texte indépendants} si la référence du produit y était écrite
quelque part. Verdict : \textbf{11 photos seulement} sont « récupérables » : la référence est
visible dans l'image, mais le modèle ne l'a pas recopiée. Un meilleur prompt pourrait donc
corriger ces 11 cas, soit \textbf{+7,6 points au maximum} sur 144 photos. Les \textbf{49 autres}
n'ont \textbf{aucune trace écrite} : aucune formulation ne peut inventer ce qui n'est pas
photographié — la cible de 70\,\% reste hors d'atteinte par le prompt. À l'inverse, la
\textbf{classification} (le type d'objet) tient bien : prise 12/12 · enceinte 76/82 · ampoule 4/5 ·
routeur 29/41 · caméra 1/4 (seulement 4 caméras : trop peu pour conclure). Le seul vrai point
faible est la confusion \emph{routeur} ↔ \emph{enceinte} (4 cas), logique pour des boîtiers
ressemblant à des enceintes connectées.

% =====================================================================
\section{Lire le texte des étiquettes (OCR) : un aide, pas une preuve}
% =====================================================================
Deux logiciels de lecture ont été comparés (PaddleOCR et OnnxTR). Deux plutôt qu'un, parce qu'on
retient un texte dès que \textbf{l'un ou l'autre} le lit (décision du 15/09) : cela fait
retrouver plus de textes réels, sans jamais changer la référence — qui reste l'étiquetage humain.
Trois mesures ont été faites, de la plus facile à la plus réaliste :
\begin{center}
\footnotesize\sffamily
\begin{tabular}{@{}p{0.24\linewidth}p{0.30\linewidth}p{0.40\linewidth}@{}}
\toprule
\textbf{Situation testée} & \textbf{Résultat} & \textbf{Ce que ça veut dire}\\
\midrule
On montre la zone du texte & 66\,\% de réussite (1\,384 textes) & Si on lui indique o\`u regarder, il recopie environ 2 textes sur 3.\\
Photo entière, gros plan & 35\,\% (600 photos, 2\,604 mots) & Il lit environ 1 mot sur 3.\\
Vraie chaîne : trouver puis lire & \textbf{10\,\%} (12 cas sur 120) & \textbf{Le chiffre honnête : 1 cas sur 10.} On ne peut pas compter sur l'OCR comme source de vérité.\\
\bottomrule
\end{tabular}
\end{center}
Dans le produit, l'OCR sert seulement à deux choses : remplir le champ « texte lu » et dire si une
photo contient du texte. Un contrôle vérifie qu'on ne triche pas avec : quand le modèle affirme
avoir lu une étiquette, on vérifie que le texte est bien présent — c'est vrai \textbf{94 fois sur
100} (75/80).

% =====================================================================
\section{La base de failles locale : ce qui est prouvé, et le choix à trancher}
% =====================================================================
\subsection{Ce qui a été importé, et pourquoi c'est traçable}
Les fichiers annuels officiels du NVD (l'agence américaine qui recense les failles) ont été
téléchargés \textbf{une fois} pour 2021–2024, conservés sur disque et \textbf{signés par empreinte}
(une empreinte est l'ADN numérique du fichier). Ainsi, un rapport pourra dire exactement quelle
version de la base il a consultée. Dans ces fichiers, chaque faille est marquée \emph{matériel}
(\code{h}), \emph{application} (\code{a}) ou \emph{système} (\code{o}) ; le projet n'a gardé que
le \emph{matériel} : \textbf{121\,719 failles connues} (une « CVE » est le nom public d'une faille,
ex. CVE-2024-1234), \textbf{324\,470 règles de correspondance} entre failles et produits, et
\textbf{873 fabricants}. Détail qui compte pour l'audit : les 3\,559 failles \textbf{retirées}
par le NVD sont conservées avec leur statut, pour qu'un rapport antérieur reste explicable.

\subsection{Comment le moteur décide}
Une fiche NVD n'est pas une simple ligne : c'est un \textbf{arbre logique} (conditions ET / OU /
EXCLUSION) dont les feuilles portent des \textbf{plages de versions}. Le moteur, écrit pour être
relu, fait donc : une première requête large en base pour éliminer l'essentiel ; puis une décision
en Python, testable ; et il affiche toujours sa \textbf{force de preuve} : produit exact > plage de
version > simple vendeur — jamais mélangées. Enfin, quand la version du logiciel interne est
inconnue (le cas normal du projet), le résultat est marqué \code{unresolved} : ni « concerné », ni
« non concerné » sans le dire. C'est la règle qui évite le pire défaut possible d'un tel outil : un
faux négatif silencieux.

\subsection{Les preuves de la porte G1 (29/09)}
\textbf{On peut relancer l'import sans risque.} Vérifié par \textbf{trois passes} sur une base
neuve, dont une avec des fichiers \emph{renommés} : c'est la seule qui réécrit vraiment les
données (la deuxième passe saute les fichiers grâce à leur empreinte). Résultat : comptes
identiques, \textbf{0 doublon}.

\textbf{Le moteur donne les bonnes réponses.} 60 couples « fabricant + produit » ont été choisis
dans toutes les catégories (beaucoup de failles, une seule, aucune…), et la réponse attendue a été
recalculée \emph{indépendamment} à partir des fichiers bruts. Le moteur donne \textbf{exactement la
même réponse 60 fois sur 60 : 0 manqué, 0 inventé}. Limites écrites noir sur blanc : la branche
« exclusion » n'est jamais déclenchée par les données réelles, et 324\,470 lignes ne sont pas
« 324\,470 failles exploitables ».

\textbf{C'est rapide.} 0,27 milliseconde à 95\,\% des cas (figure ci-dessous), pour une cible de
50 ms : \textbf{près de 200 fois plus rapide} que le minimum demandé.

\begin{center}
\begin{minipage}[t]{0.47\linewidth}
\begin{tikzpicture}
\begin{axis}[width=\linewidth,height=3.5cm,
  ybar,bar width=10pt,ymode=log,log origin=infty,ymin=0.02,ymax=100,
  symbolic x coords={p50,p95,max,cible}, xtick=data,
  ylabel={ms (échelle log)}, nodes near coords, enlarge x limits=0.25,
  grid=major, grid style={primary!12}, axis line style={primary!50}, tick style={primary!50},
  every node near coord/.append style={font=\scriptsize\sffamily\bfseries}]
\addplot[fill=ok!60,draw=ok!80] coordinates {(p50,0.062) (p95,0.267) (max,2.043)};
\addplot[fill=fail!60,draw=fail!80] coordinates {(cible,50)};
\end{axis}
\end{tikzpicture}
\captionof{figure}{Recherche de faille : médiane, 95\,\% des cas, pire cas, et la cible du contrat.
Tout tient sous 2,1 ms.}\label{fig:lat}
\end{minipage}\hfill
\begin{minipage}[t]{0.49\linewidth}
\begin{tikzpicture}
\begin{axis}[width=\linewidth,height=3.5cm,
  ybar,bar width=12pt,ymode=log,log origin=infty,ymin=1,ymax=100000,
  symbolic x coords={concernées,gardées}, xtick=data,
  ylabel={failles (échelle log)}, nodes near coords, enlarge x limits=0.35,
  grid=major, grid style={primary!12}, axis line style={primary!50}, tick style={primary!50},
  every node near coord/.append style={font=\scriptsize\sffamily\bfseries}]
\addplot[fill=fail!55,draw=fail!80] coordinates {(concernées,11714)};
\addplot[fill=ok!55,draw=ok!80] coordinates {(gardées,10)};
\end{axis}
\end{tikzpicture}
\captionof{figure}{Pour les marques de nos photos, 11\,714 failles les concernent mais \textbf{10
seulement sont gardées} par le filtre « matériel » (0,09\,\%).}\label{fig:parth}
\end{minipage}
\end{center}

\begin{atrancher}
{\footnotesize Le filtre ne conserve que les règles classées « matériel ». Conséquence mesurée le
29/09 : pour \textbf{google 2} failles gardées (1\,680 « application » et 3\,881 « système »
écartées) · \textbf{apple 2} · \textbf{huawei 1} · \textbf{samsung 3} · \textbf{netgear 2} · et
\textbf{0 pour d-link, tp-link, zyxel, amazon, philips, hikvision, mikrotik, belkin}. Autrement
dit : auditer un routeur ou une caméra de nos photos afficherait « aucune faille connue » — une
phrase que le projet s'interdit d'écrire, puisque ces failles existent. Trois sorties possibles :
\textbf{(a)} élargir l'import aux parties \emph{application} et \emph{système} ; \textbf{(b)}
déclarer le périmètre « matériel nu » et refuser explicitement au lieu de conclure ; \textbf{(c)}
élargir en conservant la force de preuve affichée. \textbf{C'est la décision la plus urgente du
projet} et elle relève du contrat, pas du code.}
\end{atrancher}

% =====================================================================
\section{Les machines de l'école : ce qui a été construit et prouvé}
% =====================================================================
Trois ordinateurs identiques (carte RTX 3060 de 12 Go) font tourner le modèle d'IA. Ils ne sont
joignables que depuis l'école : on entre par une machine de passage (le « bastion »), puis on
atteint le réseau local \code{172.26.145.x}. Le service d'IA écoute sur le port 8000 de chaque
machine, relié \emph{à l'adresse du réseau local} — joignable dans le laboratoire, invisible
d'internet. Ce n'est pas un réglage par défaut : c'est une clause du contrat de sécurité
du projet.

\begin{center}
\begin{minipage}[t]{0.38\linewidth}
\begin{tikzpicture}[node distance=3.5mm and 5mm, font=\scriptsize\sffamily]
\node[boite] (poste) {\faUser\ Poste};
\node[boite,right=of poste] (gw) {\faDoorOpen\ bastion};
\node[boiteok,right=of gw] (z22) {\faServer\ z22\code{.72}};
\node[boiteok,above=2mm of z22] (z21) {\faServer\ z21\code{.71}};
\node[boiteok,below=2mm of z22] (z23) {\faServer\ z23\code{.73}};
\draw[fleche] (poste)--(gw); \draw[fleche] (gw)--(z22); \draw[fleche] (gw)--(z21); \draw[fleche] (gw)--(z23);
\begin{scope}[on background layer]
\node[draw=primary!35,dashed,rounded corners,fit=(z21)(z22)(z23),inner sep=3pt] (lan) {};
\end{scope}
\end{tikzpicture}
\captionof{figure}{Trois nœuds derrière un bastion ; réseau local \code{172.26.145.0/24}.}
\end{minipage}\hfill
\begin{minipage}[t]{0.60\linewidth}
\begin{tikzpicture}
\begin{axis}[width=\linewidth,height=4.0cm,
  xlabel={charge demandée à chaque machine (requêtes simultanées)}, ylabel={débit (req/s)},
  xmin=6,xmax=30,ymin=15,ymax=38, grid=major, grid style={primary!12},
  axis line style={primary!50}, tick style={primary!50}]
\addplot[ok,thick,mark=*] coordinates {(8,18.0)(12,24.7)(16,30.1)(18,25.9)(24,32.1)(28,35.3)};
\addplot[primary!45,thick,dashed,mark=square*] coordinates {(8,17.4)(12,24.8)(16,30.3)(18,25.8)(24,32.0)(28,36.4)};
\draw[fail,thick,dashed] (axis cs:16,15)--(axis cs:16,38) node[pos=0.6,right,font=\tiny\sffamily,color=fail]{compromis retenu (16)};
\end{axis}
\end{tikzpicture}
\captionof{figure}{Test de charge : en vert le 28/09, en pointillé le 15/09. Les deux courbes se
superposent : le résultat est reproductible. Le débit monte jusqu'à 16 requêtes simultanées par
machine (30 req/s, latence médiane 1,57 s), passe par un creux à 18 (25,9 req/s, latence
+31\,\%), puis remonte jusqu'à 35,3 req/s à 28 requêtes — mais avec 2,27 s de latence médiane.
Le compromis retenu est \textbf{16} : au-delà, la latence augmente beaucoup plus vite que le
débit.}\label{fig:banc}
\end{minipage}
\end{center}

\noindent\textbf{Comment on sait que les machines sont saines.} Chaque nœud passe une liste de
contrôles automatiques (nom de code « NOEUD DORE » quand tout est vert) : 20 contrôles rapides
(18 passent, 2 avertissements) et 22 contrôles profonds (21 passent, 1 avertissement), sur
\textbf{les trois machines}. Les avertissements sont connus : « internet joignable » est l'état
normal du laboratoire, et le seul autre a disparu après nettoyage (voir plus bas). Deux contrôles
profonds font vraiment tourner le produit : une inférence réelle (1,49 s) et le JSON qu'elle
renvoie — \textbf{le même, caractère pour caractère, sur les trois machines} : même modèle, mêmes
poids, même réponse.

\noindent\textbf{Ce que le test de charge a appris sur le serveur lui-même.} En lisant les
compteurs internes du serveur, on voit \emph{où} passe le temps d'une requête : aux paliers 16 et
18 (huit tirages), la file d'attente est restée \textbf{nulle}, et tout le ralentissement vient
du \emph{décodage} (le moment où le modèle écrit sa réponse) — +35\,\% entre 16 et 18, retrouvé à
1\,\% près sur quatre répétitions. La file ne s'allume qu'à partir de 32 requêtes par machine,
c'est-à-dire au-delà du régime publié ici. Autre constat : la moyenne d'occupation du GPU ne
distingue pas les paliers (près de 8 points d'écart entre deux passages au même niveau ; seul le
maximum, 100\,\%, est stable) : \textbf{ce qui limite, c'est la latence}, pas la capacité du GPU.

\noindent\textbf{Le trou du redémarrage est bouché.} Avant, un simple redémarrage pouvait vider la
grappe en silence : les trois services disparaissaient, et un client ne pouvait pas distinguer une
grappe vide d'une grappe occupée. Une tâche automatique au démarrage a été ajoutée (il n'y a pas de
gestionnaire de services sur ce système) : elle attend le pilote et le réseau, puis lance le
service par le chemin habituel — \textbf{75 secondes} du début à la réponse (20 s d'attente réseau
+ 55 s de démarrage du modèle). Preuve : \textbf{trois vrais redémarrages, une machine à la fois}.
Et cette preuve a trouvé deux défauts réels dans le mécanisme (deux lancements simultanés qui se
marchaient dessus ; un lancement dans le mauvais répertoire), tous deux corrigés et re-testés.

\noindent\textbf{Internet coupé, l'IA continue.} Épreuve réalisée pour de vrai sur une machine :
avec internet bloqué (y compris le DNS du campus), \textbf{19 réponses sur 19} ont été servies, à
la vitesse normale. C'est cette même campagne qui a découvert qu'un \textbf{outil de réseau tiers}
(un accès distant installé par un autre étudiant) était actif sur \emph{deux} des trois machines
et réécrivait leurs réglages réseau. Un contrôle automatique « qui d'autre écrit sur cette
machine ? » a été ajouté, puis le logiciel tiers a été \textbf{retiré proprement} (désinstallation
réversible, sauvegarde d'état avant).

\noindent\textbf{Honnêteté sur l'instrument.} Le script de mesure a lui-même évolué en
\textbf{13 versions} : chaque défaut trouvé a été corrigé et daté. Exemple parlant : une version
comptait les requêtes \emph{en échec} comme des réussites, et leurs attentes de 120 s comme des
temps normaux — elle annonçait 93\,\% de débit en trop. Une seule version (la v4) a ses fichiers
perdus par une erreur de sauvegarde ; ses chiffres restent cités, et la règle de nommage qui évite
que ça se reproduise est écrite dans le dépôt.

% =====================================================================
\section{Ce qui est vérifié, ce qui reste, et ce qu'il faut décider}
% =====================================================================
\noindent\textbf{Vérifié le 02/10.} Les \textbf{60 tests automatiques passent} : ils protègent les
fondations (la façon de compter les objets, la comparaison des barèmes, la recherche de failles)
contre les régressions. Le contrôle de propreté du code (règles « F » et « B ») passe aussi. Le
\textbf{barème de notation par photo (\code{audit-grading}) est figé} : une copie vit dans le
dépôt, scellée par empreinte (\code{30ca55a0…}), et un test échoue si une seule ligne change —
car un barème qui bouge en silence invaliderait tous les chiffres passés. Le grand audit de
septembre a de plus recalculé \textbf{chaque figure} du dossier : le corpus se reproduit à l'identique, 4 cellules ont été
corrigées (aucun verdict changé), et \textbf{trois défauts latents} ont été trouvés et corrigés —
dont une fonction qui renvoyait toujours « tout va bien » (1,000) et la règle de consolidation des
objets, désormais dans le code au lieu de scripts jetables.

\noindent\textbf{Dette assumée.} \textbf{92 fichiers modifiés ne sont pas encore enregistrés dans Git} (le
scellement de l'état dans Git reste à faire) ; les preuves des 28–29/09 vivent sur disque, pas dans
le dépôt ; et le test de charge ne peut pas être lancé depuis ce poste d'analyse (il faut être sur
le réseau de l'école) — c'est une limite d'accès, pas un défaut du code. Volumétrie : 67 fichiers
Python (≈ 16\,000 lignes), 807 Mo sur disque dont ≈ 95\,\% de données.

\noindent Les codes \textbf{AC1} à \textbf{AC12} sont les promesses numérotées du cahier des
charges ; chaque ligne dit où on en est : \OK{} tenu · \PART{} partiel · \KO{} manqué.

\begin{center}
\begin{minipage}[t]{0.49\linewidth}
\footnotesize\sffamily
\begin{tabular}{@{}llp{0.55\linewidth}@{}}
\toprule
\textbf{\#} & & \textbf{Promesse et état}\\
\midrule
AC1 & \KO & Identifier ≥ 70\,\% des objets : \textbf{20,8\,\%} — échec, cause connue\\
AC2 & \OK & Le texte annoncé est réellement présent ≥ 90\,\% : 94\,\%\\
AC3 & \OK & Réponses toujours au bon format : 144/144\\
AC4 & \OK & Temps par photo < 10 s : 95\,\% sous 3,76 s\\
AC5 & \OK & Marque inexistante inventée < 5\,\% : 0\,\%\\
AC6 & \OK & Relancer le test donne la même chose : 144/144\\
\bottomrule
\end{tabular}
\end{minipage}\hfill
\begin{minipage}[t]{0.49\linewidth}
\footnotesize\sffamily
\begin{tabular}{@{}llp{0.55\linewidth}@{}}
\toprule
\textbf{\#} & & \textbf{Promesse et état}\\
\midrule
AC7 & \OK & L'analyse tourne internet coupé : 3/3 machines\\
AC8 & \OK & Corpus figé : 144 photos, 11 exclues tracées\\
AC9 & \PART & Base de failles rejouée + 60 cas : mesuré ; relecture humaine à faire\\
AC10 & \KO & Retrouver la bonne fiche produit : cible 80\,\% du premier coup (92\,\% dans les trois premiers choix) ; aujourd'hui 37\,\% / 56\,\% — échec\\
AC11 & \TODO & Chaîne photo → PDF scellé + tests de falsification\\
AC12 & \PART & Un test par message d'erreur (20 codes) : 0/20\\
\bottomrule
\end{tabular}
\end{minipage}
\end{center}

\noindent\textbf{Où en sont les jalons.} \textbf{G0} (socle + corpus) : fermé le 24/09.
\textbf{G1} (base de failles) : les trois critères techniques sont mesurés depuis le 29/09 — il
reste la \textbf{relecture humaine des 60 cas} et la décision de filtre ci-dessus ; échéance
19/10. \textbf{G2} (recherche produits + 200 cas) : 02/11 — aujourd'hui 37\,\% au premier choix
contre 80\,\% visés, le levier identifié est une table d'alias d'environ 50 entrées. \textbf{G3/G4}
(03–30/11) : chaîne photo → PDF, journal scellé, signature, tests de falsification.
\textbf{G5–G7} (décembre) : interface de démonstration, régénération de tous les chiffres, gel
final du 15/12.

\noindent\textbf{Prochains pas, dans l'ordre.} (1)~\emph{Enregistrer l'état dans Git} : le geste le moins
coûteux et le plus utile (92 fichiers en attente). (2)~\emph{Trancher le filtre « matériel
seulement »} et relire les 60 cas : cela ferme G1. (3)~\emph{G2} : index de recherche produits + 200 cas + table
d'alias. (4)~\emph{La confiance croisée} (accord entre photos × existence de la fiche officielle)
puis le journal scellé, la signature et le PDF.

\noindent\textbf{Deux questions attendent le tuteur.} \emph{Q3} — photographier 5–8 vrais objets
(≈ 50–60 €, une séance) pour avoir des photos « propres », ou retirer ce critère en écrivant
pourquoi. \emph{Q8} — le \textbf{FCC ID} (le numéro d'homologation imprimé sur les appareils) :
aujourd'hui zéro ligne de code et zéro donnée, mais deux de ces numéros sont déjà lus correctement
par l'OCR sur 144 photos ; recommandation : \emph{mesurer} d'abord sur ~30 photos dédiées avant de
décider, car il n'existe aucun tableau officiel reliant FCC ID et fiches produits CPE.

\vspace{1.5mm}
\begin{tcolorbox}[colback=primary!4,colframe=primary!45,boxrule=0.5pt,arc=2pt,
  left=6pt,right=6pt,top=5pt,bottom=5pt]
{\footnotesize\sffamily\textbf{En une phrase.} Au 2 octobre 2026, l'infrastructure est finie et
rejouable, la chaîne des failles est prouvée mais attend une décision de filtre, et la
reconnaissance est \emph{mesurée en échec} face à sa cible — mais son obstacle est identifié,
chiffré et actionnable : le système \textbf{lit} les marques sur les étiquettes (77\,\% de marques
justes quand le texte est lisible), il ne les \textbf{reconnaît} pas encore (10\,\% sinon). Le
prochain levier n'est pas de mieux questionner le modèle : c'est de prendre des photos où la
marque est vraiment lisible.}
\end{tcolorbox}

{\scriptsize\sffamily\color{primary!80}
\noindent\textbf{Sources :} \code{cahiers/CDC-v3.1}, \code{CDS-v2.1} · \code{docs/06} (§8–§21) ·
\code{docs/02} (décisions D1–D15) · \code{benchmarks/cve.md}, \code{results.md} ·
\code{infra/README.md} · \code{evidence/2026-09-28/}, \code{2026-09-29/} · \code{tests/}.
Chaque chiffre de ce document s'y retrouve daté.}

\end{document}
